% This appendix is for coauthor checking. Toggle it in main.tex.
\section{計算で用いた分布とパラメータ}
\label{app:distribution_defs}

本付録では, 数値計算で用いた分布のパラメータ化, 確率密度関数, および主要な計算条件をまとめる. 
ここで, \(\Gamma(\cdot)\) はガンマ関数を表す.

\subsection{確率密度関数}

\noindent\textbf{Gumbel 分布.}
\(\alpha\) を位置母数, \(\beta>0\) を尺度母数とし, \(z=(x-\alpha)/\beta\) とおく. 
このとき
\[
f(x;\alpha,\beta)
=
\frac{1}{\beta}
\exp\{-z-\exp(-z)\},
\quad -\infty<x<\infty .
\]

\noindent\textbf{一般化極値分布 (GEV).}
\(\mu\) を位置母数, \(\sigma>0\) を尺度母数, \(k\) を形状母数とし,
\[
t=1+k\frac{x-\mu}{\sigma}
\]
とおく. 
\(t>0\) の範囲で, \(k\ne0\) のとき
\[
f(x;k,\sigma,\mu)
=
\frac{1}{\sigma}
t^{-1/k-1}
\exp\{-t^{-1/k}\}
\]
である. 
\(k=0\) は Gumbel 分布への極限として扱った.

\noindent\textbf{対数ピアソンIII型分布 (LP3).}
\[
\log X = c + aY,\quad Y\sim \mathrm{Gamma}(b,1)
\]
とし, \(a>0\), \(b>0\) とする. 
\(z=(\log x-c)/a\) とおくと, \(x>\exp(c)\) に対して
\[
f(x;a,b,c)
=
\frac{z^{b-1}\exp(-z)}
{a x \Gamma(b)}
\]
である.

\noindent\textbf{平方根指数型最大値分布 (SqrtEt).}
\(a>0\), \(b>0\) とする. 
\(x\ge0\) に対して
\[
f(x;a,b)
=
\frac{ab}{2}
\exp\left[
-\sqrt{bx}
-a\left(1+\sqrt{bx}\right)\exp\{-\sqrt{bx}\}
\right]
\]
である.

\noindent\textbf{平行移動指数分布 (Exp).}
\(c\) を位置母数, \(\mu>0\) を尺度母数とする. 
\(x\ge c\) に対して
\[
f(x;c,\mu)
=
\frac{1}{\mu}
\exp\left\{-\frac{x-c}{\mu}\right\}
\]
であり, \(x<c\) では \(f(x)=0\) とする.

\noindent\textbf{3母数対数正規分布 (LN3).}
\(c\) を位置母数, \(\mu_L\) と \(\sigma_L>0\) を対数空間での位置母数および尺度母数とする. 
\[
y=\frac{\log(x-c)-\mu_L}{\sigma_L}
\]
とおくと, \(x>c\) に対して
\[
f(x;c,\mu_L,\sigma_L)
=
\frac{1}{(x-c)\sigma_L\sqrt{2\pi}}
\exp\left(-\frac{y^2}{2}\right)
\]
である.

\subsection{真の分布として用いたパラメータ}

\begin{center}
\small
\setlength{\tabcolsep}{5pt}
\begin{tabular}{lll}
\hline
分布 & パラメータ & 値 \\
\hline
Gumbel & \(\alpha,\beta\) & \(97.7438,\ 30.7042\) \\
GEV & \(k,\sigma,\mu\) & \(0.0910,\ 29.5452,\ 96.2345\) \\
LP3 & \(a,b,c\) & \(0.0683,\ 24.3314,\ 3.0352\) \\
SqrtEt & \(a,b\) & \(190.2683,\ 0.5685\) \\
Exp & \(c,\mu\) & \(54.5000,\ 61.6320\) \\
LN3 & \(c,\mu_L,\sigma_L\) & \(31.3965,\ 4.3254,\ 0.4805\) \\
\hline
\end{tabular}
\end{center}

\subsection{主な計算条件}

\begin{center}
\small
\setlength{\tabcolsep}{5pt}
\begin{tabular}{ll}
\hline
項目 & 設定 \\
\hline
標本サイズ & \(N=50,100,150\) \\
反復数 & \(100\) \\
乱数シード & \(42\) \\
プロッティングポジション & Cunnane 式 \\
SLSC 閾値 & \(0.04\) \\
ジャックナイフ幅の参照再現期間 & \(T_{\mathrm{ref}}=100\) \\
\hline
\end{tabular}
\end{center}
